\section{Recopie de texte} 
%\label{sec:5:1:recopie} 

Dans le programme {\idt CroixGrecque\_1} ci-dessous on retrouve quatre fois la même
séquence de commandes pour dessiner les quatre
bran\-ches identiques de la croix et tourner. Chaque branche se décompose en deux cotés égaux
de lon\-gueur 3 et une extrémité de lon\-gueur 1 (l'unité
étant la longueur d'un côté d'un carreau  du quadrillage).

\vspace{\baselineskip}
\vspace*{3cm}%
\includevspacePostScript[-\baselineskip]{CroixGrecque}{0.95}{0.95}
\vspace*{-3cm}%
\label{prog:CroixGrecque_1}
\input python/CroixGrecque_1.tex

Cette situation où l'on trouve la même séquence de commandes du \robot\ 
plusieurs fois de suite dans un programme est très fréquente car les 
dessins géométriques sont souvent composés de morceaux identiques. Il est alors
très fastidieux de taper plusieurs fois le même texte. On peut l'éviter
en utilisant les possibilités offertes par l'ordinateur de copier un texte
déjà enregistré dans sa mémoire et de le recopier ailleurs dans le 
programme. La manière de procéder est expliquée pour les 
PC (sous Windows ou Linux) à la figure \ref{fig:PC-copier-coller} 
(page \pageref{fig:PC-copier-coller})
et pour le Macintosh à la figure \ref{fig:mac-copier-coller} 
(page \pageref{fig:mac-copier-coller}).
\input COPIER-COLLER-PC.tex
\input COPIER-COLLER-MAC.tex

Faire beaucoup de recopies à la main est long. De plus on imagine bien que 
dans un programme de plusieurs pages, le lecteur a du mal à retrouver les 
parties de texte identiques. Ce phénomème est illustré dans le
programme {\idt Nombre888\_1} ci-dessous. La difficulté serait
évidemment encore plus grande si les parties de texte identiques se
trouvaient sur des pages différentes.

\vspace{\baselineskip}
\vspace*{2cm}%
\includevspacePostScript[-\baselineskip]{Nombre_888}{0.95}{0.95}
\vspace*{-2cm}%
\input python/Nombre888_1.tex

\section{Exemples de procédures} 
%\label{sec:5:2:ex-procedures} 

Plutôt que de faire des recopies de textes à la main soi-même, on peut 
demander à l'interpréteur \python\ de le faire à notre place en utilisant 
des {\em procédures}\index{Procédure}. Commençons par réécrire le programme 
{\idt CroixGrecque\_1} (page \pageref{prog:CroixGrecque_1}) en uti\-li\-sant des procédures pour mettre en 
évidence la décomposition de la croix grecque en 
quatre bran\-ches identiques, chacune d'elles se décomposant en deux 
cotés égaux de lon\-gueur 3 et une extrémité de lon\-gueur 1.

\newpage
\vspace*{3cm}%
\includePostScript{CroixGrecque}{0.95}{0.95}
\vspace*{-3cm}%
\input python/CroixGrecque_2.tex

\noindent L'emploi d'une procédure comprend~:
\begin{list}{$\scriptscriptstyle\bullet$}{
}
\item la {\em déclaration de procédure}\index{Déclaration de procédure} qui définit une séquence 
d'instructions (appelée {\em corps de la procédure})\index{Corps de procédure} et lui donne un 
nom (appelé {\em nom de la procédure})\index{Identificateur}$\:$;
\item l'{\em appel de procédure}\index{Appel de procédure} qui indique que la séquence 
d'instructions cons\-ti\-tuant le corps de la procédure doit être recopiée 
à la place du nom de la procédure.
\end{list}

\noindent Pour le lecteur du programme, l'utilisation de procédures est 
préférable aux recopies de textes. A chaque fois qu'il voit le
nom d'une procédure, il comprend qu'il s'agit de la copie d'un texte
qu'il peut consulter dans la déclaration de procédure. S'il rencontre
plusieurs appels d'une même procédure, il n'a pas à faire d'effort 
pour comprendre immédiatement qu'il s'agit du même texte qui est
recopié. Le programme est donc plus clair. N'est-ce pas le cas du
programme {\idt Nombre888\_2} ci-dessous qui est plus facilement 
compréhensible que le programme {\idt Nombre888\_1}$\:$? 

\vspace{\baselineskip}
\vspace*{2cm}%
\includevspacePostScript[-\baselineskip]{Nombre_888}{0.95}{0.95}
\vspace*{-2cm}%
\input python/Nombre888_2.tex

\section{Déclaration de procédures} 
%\label{sec:5:3:decl-procedures} 

Les déclarations des procédures\index{Déclaration de procédure} d'un programme se placent entre 
la ligne `{\mcl from}\ {\idt LibRobot}\ {\mcl import}\ {\idt *}' et les instructions du programme. 
La structure d'un programme \python\ est donc la suivante~:

\begin{center}
\begin{tabular}{@{}|p{12.5cm}|@{}} \hline
{\em Programme \python\ :}\\ \hline \hline
\begin{minipage}[t]{12.5cm}
{\medskip
\mbox{\blc\ \ \ {\mcl from}\hspace{0.3em}\blanc\hspace{0.3em}{\idt LibRobot}\hspace{0.3em}\blanc\hspace{0.3em}{\mcl import}\hspace{0.3em}\blanc\hspace{0.3em}{\idt *}}\par 
\mbox{\blc\ \ \ \cmt \#\ \small\rm Commentaire expliquant ce que fait le programme}\par 
\mbox{\blc\ \ \ \small\rm  \ldots\ Déclarations des procédures du programme \ldots}\par 
\mbox{\blc\ \ \ \small\rm  \ldots\ Instructions du programme%
\footnote{\ Y compris les appels des procédures du programme.} (non indentées) \ldots}\par 
\medskip}
\end{minipage}\\ \hline 
\end{tabular}
\end{center}


\noindent En \python, une déclaration de procédure a la forme suivante
(nous utiliserons des formes plus élaborées dans la suite)~:

\begin{center}
\begin{tabular}{@{}|p{12.5cm}|@{}} \hline
{\em Déclaration de procédure~:}\\ \hline \hline
\begin{minipage}[t]{12.5cm}
{\medskip
\mbox{\blc\ \ \ \cmt \#\ \small\rm Spécification : ce que fait la procédure}\par 
\mbox{\blc\ \ \ {\mcl def}\hspace{0.3em}\blanc\hspace{0.3em}\idt nomDeLaProcedure():}\par 
\mbox{\blc\ \ \ \blk\blk\blk\blk\small\rm \ldots\ Corps de la procédure%
\footnote{\ Y compris les appels de procédures précédemment déclarées.}
 ({\bf obligatoirement indenté}) \ldots}\par 
\medskip}
\end{minipage}\\ \hline 
\end{tabular}
\end{center}

Le mot anglais {\mcl def} est l'abréviation de {\em define} (définir).
Le {\idt  nomDeLaProcedure} est choisi par le programmeur
à sa convenance, de façon à évoquer ce
à quoi sert cette procédure. Les parenthèses {\idt ()} qui le suivent indiquent
que c'est une procédure, qu'il faudra appeler en écrivant {\idt nomDeLaProcedure()}, 
exactement comme les commandes du \robot\ {\idt av()}, {\idt pqtd()}, etc. 
D'ailleurs les commandes du \robot\ sont des procédures déclarées dans la 
bibliothèque {\idt LibRobot}. La ligne {\mcl def} se termine obligatoirement par 
deux points `{\idt :}' qui annoncent le corps de la procédure.

\index{Indentation}Les instructions du corps de la procédure sont écrites sur les lignes 
suivantes, toutes décalées vers la droite d'un même nombre de blancs, 
habituellement 4. On dit qu'elles sont {\em indentées}. 
Comme dans le programme, on peut écrire une ou plusieurs instructions par ligne, séparées 
par des points-virgules. C'est l'indentation, et elle seule, qui indique à \python\ 
quelles instructions font partie du corps de la procédure~: le corps se termine à 
la première ligne qui n'est plus indentée. Les blancs en début de ligne ne sont donc 
pas une simple question de présentation~: ils font partie de la syntaxe de \python. 
Si l'on oublie d'indenter la première ligne du corps, \python\ signale par exemple
`{\tt IndentationError: expected an indented block after function definition on line 3}' 
(erreur d'indentation~: un bloc indenté est attendu après la définition de fonction de 
la ligne 3). Si au contraire on indente par erreur une instruction du programme, il signale 
`{\tt IndentationError: unexpected indent}' (indentation inattendue). Et si une ligne du 
corps n'a pas le même nombre de blancs que les précédentes, il signale 
`{\tt IndentationError: unindent does not match any outer indentation level}' 
(le retrait ne correspond à aucun niveau d'indentation). Une ligne vide après le corps 
de chaque procédure, comme dans le programme {\idt Nombre888\_2} ci-dessus, 
n'est pas obligatoire mais rend bien visible la fin de la procédure.

{\bf Attention~:} le corps d'une procédure doit {\bf impérativement} être indenté, et 
{\bf toutes} ses lignes du même nombre de blancs. Une ligne du corps que l'on oublie 
d'indenter n'est pas toujours signalée comme une erreur~: elle ne fait simplement plus 
partie de la procédure. Dans l'exemple ci-dessous, {\idt pqtd()} n'est pas indentée~:
\python\ l'exécute une seule fois, au début du programme, et non à chaque appel de 
{\idt coin()}. Au lieu d'un carré, le \robot\ trace un trait droit de huit pas, sans 
aucun message d'erreur~:

{\parindent \parindentprog
\medskip
\noindent\begin{minipage}{\linewidth}\parindent \parindentprog
\vspace{\neglskip}\mbox{\blc {\mcl def}{\idt \ coin():}}\par
\vspace{\neglskip}\mbox{\blc \blk\blk{\idt av();\ av()}}\par
\vspace{\neglskip}\mbox{\blc {\idt pqtd()}\ \ \ \ \ \ \ \ \ \ {\cmt \#\ Erreur\ :\ pas\ indentée,\ hors\ de\ la\ procédure}}\par
\vspace{0.3\baselineskip}
\vspace{\neglskip}\mbox{\blc {\idt coin();\ coin();\ coin();\ coin()}}\par

\end{minipage}
\medskip }

\noindent Pour que le \robot\ tourne à chaque coin, il faut indenter {\idt pqtd()} 
comme {\idt av();\ av()}, avec les mêmes quatre blancs.

\medskip
\noindent {\it Remarque~:} en \python, les procédures sont des cas particuliers de 
ce que l'on appelle des {\em fonctions}\index{Fonction}. Nous en reparlerons 
plus tard.


\section{Identificateurs} 
%\label{sec:7:2:Ident} 
Le nom d'une procédure s'appelle un 
{\em identificateur}\index{Identificateur} en \python. En voici 
une définition précise à laquelle il
faudra se référer à chaque fois que nous parlerons d'iden\-ti\-fi\-ca\-teurs~:

\begin{center}
\begin{tabular}{@{}|p{12.5cm}|@{}} \hline
{\em Identificateur~:}\\ \hline \hline
\begin{minipage}[t]{12.5cm}
{\medskip
Un {\em identificateur} est une séquence de lettres 
minuscules ou majuscules non accentuées, de chiffres ou
de soulignés `\_' qui ne commence pas par un chiffre et qui est différente des
mots anglais réservés par \python\ comme {\mcl from}, {\mcl import}, {\mcl def}, {\mcl for}, 
{\mcl in}, {\mcl if}, {\mcl while}, {\mcl return}, {\mcl True}, etc.
\medskip}
\end{minipage}\\ \hline 
\end{tabular}
\end{center}

Par exemple {\idt cote}, {\idt croixGrecque}, {\idt croix\_grecque} ou {\idt cote2} sont des 
identificateurs alors que {\idt 2cote}, {\idt cote-gauche} ou {\idt for} n'en sont pas. 
Si l'on écrit `{\mcl def}\ {\idt 2cote():}', \python\ répond 
`{\tt SyntaxError: invalid decimal literal}' (nombre décimal incorrect) car il croit 
lire un nombre$\:$; pour `{\mcl def}\ {\idt cote-gauche():}' ou `{\mcl def}\ {\idt for():}', il 
répond simplement `{\tt SyntaxError: invalid syntax}' (syntaxe incorrecte).
Les majuscules et les minuscules sont distinguées~: {\idt Cote} et {\idt cote} sont deux 
identificateurs différents. Si l'on déclare `{\mcl def}\ {\idt Cote():}' puis que l'on 
appelle {\idt cote()}, \python\ répond 
`{\tt NameError: name \textquotesingle{}cote\textquotesingle{} is not defined. Did you mean: \textquotesingle{}Cote\textquotesingle{}?}' 
(le nom {\idt cote} n'est pas défini. Vouliez-vous dire {\idt Cote}$\:$?).
\python\ permet de commencer un identificateur par une majuscule, mais l'usage est de 
commencer les noms de procédures par une minuscule. Pour rendre 
lisible un identificateur formé de plusieurs mots, on
met une majuscule au début de chaque mot sauf le premier, comme dans 
{\idt croixDeSaintAndre} ou {\idt bordInterieurDuVerre}\footnote{\ Beaucoup de programmeurs 
\python\ préfèrent séparer les mots par des soulignés, comme dans {\idt croix\_de\_saint\_andre}.}.

\section{Appel de procédures} 
%\label{sec:5:4:app-procedures} 

\index{Appel de procédure}Pour appeler une procédure, il suffit d'écrire son nom suivi 
de {\idt ()}. L'appel {\idt nomDeLaProcedure()}
est alors remplacé par le corps de la procédure\index{Corps de procédure} c'est-à-dire par la liste
d'instructions indentées qui suit la ligne `{\mcl def}\ {\idt nomDeLaProcedure():}' dans sa déclaration.

On notera que si l'on déclare une procédure sans jamais l'appeler, elle 
ne sert absolument à rien. 

On peut appeler une procédure dans le programme ou dans une autre procédure.
Par exemple dans le programme {\idt CroixGrecque\_2}, on appelle 
la procédure {\idt croix} qui appelle la procédure {\idt branche} qui elle-même
appelle les procédures {\idt cote} et {\idt extremite}. 

Il faut impérativement respecter la règle qu'une procédure doit 
être
dé\-cla\-rée avant d'être ap\-pe\-lée. Supposons par exemple que l'on ait
déclaré une pro\-cé\-du\-re sous le nom {\idt nomProc}. Supposons également
que par suite d'une faute de frappe, on ait appelé cette procédure
sous le nom {\idt nomProg}. Quand l'exécution arrive à cet appel, l'interpréteur 
s'arrête en indiquant l'erreur
`{\tt NameError: name \textquotesingle{}nomProg\textquotesingle{} is not defined}' (le nom {\idt nomProg} n'est pas défini). 
Ne sachant pas que 
nous avons écrit par erreur {\idt nomProg} au lieu de {\idt nomProc}, l'interpréteur 
suppose que la procédure {\idt nomProc} a été déclarée mais n'a jamais été utilisée.
C'est inutile mais autorisé. Il constate également que la procédure
{\idt nomProg} est appelée mais n'a pas été déclarée. C'est interdit.
Il remarque cependant que les deux noms se ressemblent et suggère
`{\tt Did you mean: \textquotesingle{}nomProc\textquotesingle{}?}' (vouliez-vous dire {\idt nomProc}$\:$?).%
\footnote{\ Nous parlons souvent des programmes de manière 
anthropomorphique, c'est-à-dire en leur attribuant des caractères 
propres à l'homme. En fait nous pensons aux programmeurs qui les
ont écrits. Par exemple, l'ordinateur ne se trompe pas$\:$; c'est son 
constructeur qui s'est trompé ou le programmeur qui a fait une erreur 
logique dans son programme.} 
On obtient la même erreur `{\tt NameError}' si le programme appelle une procédure
dont la déclaration est placée plus loin~: \python\ exécute les déclarations dans 
l'ordre où elles sont écrites, et au moment de l'appel il ne connaît pas encore la procédure. 
Comme toutes les erreurs découvertes pendant l'exécution, ces erreurs n'apparaissent 
qu'une fois que le \robot\ a fait la partie du dessin qui précède l'appel fautif.

\section{Commandes du robot définies par des procédures} 
%\label{sec:5:5:robot-cmd-procs}

Certaines commandes du \robot\ sont équivalentes à une séquence
de commandes élémentaires et peuvent donc s'exprimer avec des procédures.
Par exemple, pour pivoter d'un quart de tour, il suffit de pivoter deux fois
de suite d'un huitième de tour. On a donc~:
\input python/Pivotements1.tex
\noindent Voici de nouvelles commandes du \robot\ qui peuvent s'exprimer à l'aide de
procédures en utilisant des commandes que nous connaissons déjà~:
\input python/Pivotements2.tex
\noindent\rule{137mm}{0.1mm} %-------------------------------------------------------
\begin{exercices}{
Écrire des programmes \python\ comportant des pro\-cé\-du\-res pour 
réaliser les croix,
%de la figure \ref{fig:Croix}, 
l'étoile astroïdale, la loupe ou le double carré cranté
%de la figure \ref{fig:EtLoCa} 
ci-dessous.
\noindent $\Box$}\end{exercices}
\begin{figure}[htbp]
 \vspace{-\baselineskip}
 \includevspacePostScript[1.7cm]{Croix}{0.95}{0.95}
 \caption{Croix potencée, double, papale, tréflée et de Saint André}
 %\label{fig:Croix}
\end{figure}
\begin{figure}[htbp]
 \includevspacePostScript[5.1cm]{EtoileLoupeCarre}{0.95}{0.95}
 \caption{Étoile astroïdale, loupe et double carré cranté}
 %\label{fig:EtLoCa}
\end{figure}

\clearpage

\begin{corrige}{
Nous donnons un exemple de programme, les
autres se trouvant corrigés dans le dossier {\idt livre/\allowbreak 05\_procedures}.
\input python/EtoileAstroidale.tex
\noindent $\Box$}\end{corrige}

\vfill

%%%%%%%%%%%%%%%%%%%%%%%%%%%%%%%%%%%%%%%%%%%%%%%%%%%%%%%%%%%%%%%%%%%%%%%%%%%
\noindent \framebox[\textwidth]{
\begin{minipage}{\myparboxwidth}
\setlength{\parindent}{2em}
\small \sf
\vspace{\abovedisplayskip}
\section*{Addenda}
La notion de procédure\index{Procédure} introduite par Von Neumann et Turing
fût utilisée en 1949
sous le nom de {\em routine} sur l'EDSAC ({\em Electronic Delay
Storage Automatic Computer}) généralement considéré 
comme le premier ordinateur à programme enregistré ayant 
fonctionné. Cette machine fût construite dès 1947 à l'université
de Cambridge en Angleterre par l'équipe de Wilkes\footnote{\ Maurice
Wilkes, l'un des inventeurs du radar.} et exécuta ses premiers programmes
en Mai 1949 \cite{levy1}.

Les procédures illustrent le concept d'{\em abstraction}\index{Abstraction}~: les 
propriétés d'un objet sont isolées afin de
les utiliser ou de les étudier à part. Plus précisément, une 
procédure donne un nom à un algorithme de construction d'une partie
d'un dessin. Dans le corps de la procédure, 
on s'intéresse à la façon de construire le dessin. Dans les appels 
on s'intéresse uniquement au dessin terminé, tout en désirant complètement ignorer comment 
le dessin est obtenu. On sépare donc l'objet et la façon de le 
construire, le quoi du comment. 

\vspace{\abovedisplayskip}
\end{minipage}
}
%%%%%%%%%%%%%%%%%%%%%%%%%%%%%%%%%%%%%%%%%%%%%%%%%%%%%%%%%%%%%%%%%%%%%%%%%%%%
%%%%%%%%%%%%%%%%%%%%%%%%%%%%%%%%%%%%%%%%%%%%%%%%%%%%%%%%%%%%%%%%%%%%%%%%%%%
\noindent \framebox[\textwidth]{
\begin{minipage}{\myparboxwidth}
\setlength{\parindent}{2em}
\small \sf
\vspace{\abovedisplayskip}
Nous trouvons utile de décrire le dessin réalisé par une procédure
dans un commentaire 
placé au début de la procédure que l'on appelle la
 {\em spécification}\index{Spécification} de la procédure. On remarquera que pour qu'une procédure 
puisse être utilisée sans jamais consulter son corps, il faut que la 
spécification de la procédure décrive quel dessin elle fait mais 
également quelles sont les positions et orientations
du \robot\ avant et après la construction du dessin. Il se dégage ainsi
une notion d'{\em invariant}\index{Invariant} sur laquelle nous reviendrons dans les chapitres
suivants sur l'itération. 

Les procédures sont à la base des méthodes 
de structuration des très grands programmes~: pour être 
compréhensibles, les programmes de plusieurs milliers de lignes 
doivent être découpés en procédures ne dépassant pas 
quelques pages. Les procédures constituent donc un moyen élégant pour 
segmenter logiquement un programme. Cependant le concept de procédure
est insuffisant pour structurer des programmes de plusieurs millions de
lignes. On utilise alors la notion de {\em module} (on dit également {\em unité})
et de {\em bibliothèque} pour regrouper des ensembles de procédures.

Des procédures regroupées dans une
{\em bibliothèque} (en \python\ on dit aussi {\em module}) peuvent être 
fournies aux programmeurs pour les aider dans leurs tâches. Pour utiliser 
une procédure, le programmeur doit connaître son nom et sa 
spécification sans avoir besoin
de connaître son fonctionnement interne. La ligne `{\mcl from}\ {\idt LibRobot}\ {\mcl import}\ {\idt *}' au début
du programme a pour but d'indiquer à l'interpréteur quelle est la 
bibliothèque utilisée dans ce programme. Les commmandes du \robot\ sont
des procédures regroupées dans des bibliothèques.

L'analyse de dessins géométriques pour en dégager une structure que
l'on peut décrire algorithmiquement à l'aide de procédures est très
formatrice pour l'esprit. En ce sens l'apprentissage de la programmation
peut jouer pour les enfants un rôle formateur analogue à celui que 
jouait autrefois la géométrie euclidienne.

\vspace{\abovedisplayskip}
\end{minipage}
}
%%%%%%%%%%%%%%%%%%%%%%%%%%%%%%%%%%%%%%%%%%%%%%%%%%%%%%%%%%%%%%%%%%%%%%%%%%%%
\newpage
\thispagestyle{empty}